\renewcommand{\footerfilename}{CriticalConcepts.tex}

\section{Critical Concepts}


\dhl{DO NOT EVER} use unicode. Anywhere. IRAF rendering uses specialized ASCII only fonts.\\
\dhl{DO NOT EVER} use spaces in filenames. Professional reduction software chokes on this. \\
\dhl{DO NOT EVER} use plus, minus, or non-ASCII (fancy UNICODE chars) in filenames, FITS headers
etc.

The \dhl{date} of an observation's dataset is specifed as the calendar date
of sunset at the observatory in the observatory's time zone. The use
of UT really messes things up.\
\dhl{DO KEEP FILENAMES SHORT} The filename length essentially hardwired into IRAF is
around 23 ASCII characters. Longer than this length messes up a lot of column headers.

\dhl{DO DO DO PUT TARGET AND REFERENCE} info into the filename. You
can then update the headers when the night's dataset is being setup
for analysis.\
\dhl{DO NOT USE} special characters or spaces, use underscore for
space, \dhl{m} for minus \dhl{p} plus etc. If the software adds UT
timestamps (FITS keyword DATE-OBS) to files, use them\footnote{Zip files in general and
Windows does not preserve systems' dates/times with filenames reliably.}

Make a directory called
\dhl{{\textasciitilde{}/Desktop/Today}}. Hardwire this path into all
the software that writes files somewhere during the observation. Thus,
you do not have to visit those packages (and risk forgetting one) each
night. At the end of the night rename to taste, move into a collection
director; make a new directory and name it
\dhl{{\textasciitilde{}/Desktop/Today}}. A rinse/repeat
thing.\ \dhl{DO MAKE BACKUPS} The rule of backups is 3-2-1. This
strategy means to make three copies, on at least two different types
of media and store at least one to an off-site location. It is wise to
make copies using different archive means: like tar and zip and other
methods.

Include/export as many configuration files as you can to your dataset.
Script these things.

Get as much relevant information as possible into each file's header using
the software. Develop scripts to add other details like: Latitude/Longitude,
site name, observers, telescope details etc. 


\dhl{DO} keep filename length to a minimum.

\dhl{DO DO DO} Have each observer keep a paper logbook, using an ink that
does not run when wet.
